%% chapters/04-citations-references.tex
\chapter{Citations and Cross-References}\label{chap:citations}

This chapter explains how to refer to other parts of the book and to external
sources.

\section{Cross-References}\label{sec:crossrefs}

Give every chapter, section, figure, table, equation and listing a
\cmd{label}, then refer to it with \cmd{cref}. The \pkg{cleveref} package
writes the kind of object for you, so you never type \enquote{Figure} or
\enquote{Table} by hand and the text stays correct when things are
renumbered:

\begin{lstlisting}[language={[LaTeX]TeX}]
As \cref{fig:sample} shows, ...
See \cref{sec:crossrefs} and \cref{tab:environments}.
\Cref{eq:quadratic} is the quadratic formula.
\end{lstlisting}

\noindent For example: as \cref{fig:sample} shows, a plot can be referenced
from anywhere; see \cref{sec:crossrefs} and \cref{tab:environments}.
\Cref{eq:quadratic} is the quadratic formula. Use \cmd{Cref} with a capital C
at the beginning of a sentence. Several references can be combined, as in
\cref{sec:figures,sec:tables,sec:code}.

Label names are free text, but a prefix keeps them organised:

\begin{center}
  \begin{tabular}{@{}ll@{\hspace{2em}}ll@{}}
    \toprule
    Prefix & Object & Prefix & Object \\
    \midrule
    \texttt{chap:} & chapter  & \texttt{tab:} & table \\
    \texttt{sec:}  & section  & \texttt{eq:}  & equation \\
    \texttt{fig:}  & figure   & \texttt{lst:} & listing \\
    \bottomrule
  \end{tabular}
\end{center}

\begin{note}
Place \cmd{label} \emph{after} \cmd{caption} inside floats, or right after the
sectioning command. Otherwise it may point to the wrong number.
\end{note}

\section{Citations}\label{sec:citations}

Bibliographic data lives in \file{references.bib}, one entry per source. The
key of each entry identifies it in the text. Two commands are used most:

\begin{itemize}
  \item \cmd{citep}\texttt{\{key\}} for a parenthetical citation, as in
        \enquote{\ldots\ is well documented \citep{lamport1994latex}}.
  \item \cmd{citet}\texttt{\{key\}} for a citation that is part of the
        sentence, as in \enquote{\citet{knuth1984texbook} describes the
        typesetting engine}.
\end{itemize}

\noindent Several keys can be cited together with commas:
\citep{lamport1994latex,mittelbach2004companion}. Web resources are handled in
the same way; a good place to start is the \citet{overleaf} tutorial.

A minimal entry looks like this:

\begin{lstlisting}[language={}]
@book{knuth1984texbook,
  author    = {Knuth, Donald E.},
  title     = {The {\TeX}book},
  publisher = {Addison-Wesley},
  year      = {1984}
}
\end{lstlisting}

\section{The Bibliography}\label{sec:bibliography}

The bibliography is printed at the end of the book by a single line in
\file{main.tex}:

\begin{lstlisting}[language={[LaTeX]TeX}]
\bibliography{references}
\end{lstlisting}

\noindent Citations are numeric by default. To use the author--year style,
load the style package with an option in \file{main.tex}:

\begin{lstlisting}[language={[LaTeX]TeX}]
\usepackage[authoryear]{style/booktemplate}
\end{lstlisting}

\begin{tip}
If a citation appears as \texttt{[?]}, the bibliography has not been
processed yet. On Overleaf, recompile; on your computer, run
\texttt{latexmk main.tex}, which runs BibTeX as needed.
\end{tip}

\section{Hyperlinks}\label{sec:hyperlinks}

Every cross-reference, citation and table-of-contents entry is a clickable
link. External addresses are written with \cmd{url} or, with a custom text,
\cmd{href}: \url{https://ctan.org} or \href{https://ctan.org}{the CTAN
archive}. Link colours follow the accent colour defined in
\file{style/booktemplate.sty}.
